\section*{Chapitre \ref{ch:b3:advanced-geochronology} --- Géochronologie avancée}

\begin{solution}{exo:b3:advanced-geochronology:1}
$t_{206} = \ln(1.04500)/\lambda_{238} = \qty{283.8}{\Ma}$; $t_{207} =
\ln(1.3300)/\lambda_{235} = \qty{289.6}{\Ma}$. Les âges diffèrent de
\qty{6}{\Ma}~: le zircon est discordant. Son $t_{207}$ est le plus ancien, il se place
sous la \omterm{def:b3:advanced-geochronology:concordia}{concordia}, du côté de la perte de plomb.
\end{solution}

\begin{solution}{exo:b3:advanced-geochronology:2}
$J = (e^{5.543\times10^{-10}\times2.8201\times10^7} - 1)/1.3129 =
0.015754/1.3129 = 0.01200$. Puis $t = \ln(1 + 0.01200\times10.000)/\lambda =
\ln 1.12/(5.543\times10^{-10}) = \qty{204.5}{\Ma}$.
\end{solution}

\begin{solution}{exo:b3:advanced-geochronology:3}
$T_{\mathrm{CHUR}} = \ln\bigl(1 + (0.51180 - 0.512638)/(0.1200 -
0.1966)\bigr)/\lambda = \ln(1.010940)/\lambda = \qty{1.67}{\Ga}$;
$T_{\mathrm{DM}} = \ln\bigl(1 + (0.51180 - 0.51315)/(0.1200 - 0.2137)\bigr)/\lambda
= \ln(1.014408)/\lambda = \qty{2.19}{\Ga}$.
\end{solution}

\begin{solution}{exo:b3:advanced-geochronology:4}
Sans décroissance, $t = N/P_0 = \num{2.0e5}/4.96 = \qty{40.3}{ka}$. Avec décroissance,
$t = -\ln(1 - N\lambda/P_0)/\lambda = -\ln(1 - 0.02014)/\lambda =
\qty{40.7}{ka}$, soit \qty{1}{\percent} de plus.
\end{solution}

\begin{solution}{exo:b3:advanced-geochronology:5}
Pour les petits $t$, $x \approx \lambda_{235}t$ et $y \approx \lambda_{238}t$, donc
$y \approx (\lambda_{238}/\lambda_{235})x$, pente $0.1575$. Si le plomb a été
perdu très récemment, $t_2 \approx 0$ et l'\omterm{def:b3:advanced-geochronology:concordia}{intercept inférieur} est $C(0)$,
soit l'origine. Sur une droite passant par l'origine, $y/x$ est constant, de même que le rapport
$^{207}\mathrm{Pb}^*/^{206}\mathrm{Pb}^* = x/(137.8\,y)$~: chaque cristal donne
l'âge plomb--plomb de l'\omterm{def:b3:advanced-geochronology:concordia}{intercept supérieur}, soit l'âge de formation.
\end{solution}

\begin{solution}{exo:b3:advanced-geochronology:6}
Poids $1.5625$, $1$, $1.2346$, $2.0408$, $W = 5.838$~; $\bar x =
\qty{296.09}{\Ma}$, $\sigma = 1/\sqrt W = \qty{0.41}{\Ma}$. Les écarts pondérés
au carré ont pour somme $2.32$~; avec $f = 3$, $\mathrm{MSWD} = 0.77$, inférieur à
$1 + 2\sqrt{2/3} = 2.63$~: les âges sont compatibles.
\end{solution}

\begin{solution}{exo:b3:advanced-geochronology:7}
$\lambda\zeta g\rho_d\rho_s/\rho_i = 1.55125\times10^{-10}\times350\times0.5
\times\num{1.10e6}\times0.5 = 0.014931$; $t = \ln(1.014931)/\lambda =
\qty{95.5}{\Ma}$. Si la roche s'est refroidie vite, elle date son passage dans la
zone de recuit partiel de l'apatite, vers 60 à \qty{120}{\celsius}~; si elle
s'est refroidie lentement, c'est un âge mixte, et les longueurs de traces doivent trancher.
\end{solution}

\begin{solution}{exo:b3:advanced-geochronology:8}
Par gramme~: uranium $20\times10^{-6}/238.05\times N_A = \num{5.06e16}$ atomes, dont
\num{5.02e16} d'uranium-238 et \num{3.64e14} d'uranium-235~; thorium
\num{1.038e17}, $\lambda_{232} = \ln 2/\num{1.40e10} = \qty{4.95e-11}{an^{-1}}$~;
hélium $10^{-9}N_A = \num{6.02e14}$. Le \omterm{def:b3:advanced-geochronology:cosmogenic}{taux de production} est
$8\lambda_{238}{}^{238}\mathrm U + 7\lambda_{235}{}^{235}\mathrm U +
6\lambda_{232}\mathrm{Th} = (6.23 + 0.25 + 3.08)\times10^7 =
\qty{9.57e7}{atomes.g^{-1}.an^{-1}}$, et $t = \num{6.02e14}/\num{9.57e7} =
\qty{6.29}{Ma}$. Corrigé~: $6.29/0.75 = \qty{8.39}{\Ma}$.
\end{solution}

\begin{solution}{exo:b3:advanced-geochronology:9}
$P_0/\lambda = 4.96/\num{4.994e-7} = \qty{9.93e6}{atomes.g^{-1}}$, atteinte à
\qty{90}{\percent} après $\ln 10/\lambda = \qty{4.6}{Ma}$. Pour
\qty{10}{m.Ma^{-1}} $= \qty{1e-3}{cm.an^{-1}}$~: $\rho\varepsilon/\Lambda =
2.7\times10^{-3}/160 = \qty{1.69e-5}{an^{-1}}$, $\mu = \qty{1.74e-5}{an^{-1}}$
et $N = 4.96/\num{1.74e-5} = \qty{2.85e5}{atomes.g^{-1}}$, 35 fois moins.
\end{solution}

\begin{solution}{exo:b3:advanced-geochronology:10}
Avec $W = w_1 + w_2$, $x_1 - \bar x = w_2(x_1 - x_2)/W$ et $x_2 - \bar x =
-w_1(x_1 - x_2)/W$, donc $\sum w_i(x_i - \bar x)^2 = (x_1 - x_2)^2(w_1w_2^2 +
w_2w_1^2)/W^2 = (x_1 - x_2)^2w_1w_2/W = (x_1 - x_2)^2/(\sigma_1^2 +
\sigma_2^2)$~; avec $f = 1$, c'est le MSWD. Si les deux mesurent le même âge,
$x_1 - x_2$ a pour moyenne 0 et pour variance $\sigma_1^2 + \sigma_2^2$, donc
l'espérance de son carré, divisée par cette variance, vaut 1.
\end{solution}

\begin{solution}{exo:b3:advanced-geochronology:11}
$C(2700) = (13.283, 0.52019)$, $C(500) = (0.63628, 0.08065)$. $f = (0.52019 -
0.300)/(0.52019 - 0.08065) = 0.501$, soit la moitié de son plomb. $x = 0.499\times13.283 +
0.501\times0.63628 = 6.948$. $t_{206} = \qty{1691}{\Ma}$, $t_{207} =
\qty{2105}{\Ma}$; le rapport $x/(137.8\,y) = 0.1681$ donne par itération un
âge plomb--plomb de \qty{2537}{\Ma}. Aucun des trois n'a de sens~: ce sont
des mélanges des deux événements, et seuls les intercepts les datent. L'âge
plomb--plomb constitue un minimum pour la cristallisation.
\end{solution}

\begin{solution}{exo:b3:advanced-geochronology:12}
Sans érosion~: $t = -\ln(1 - \num{1.0e6}\times\num{4.994e-7}/4.96)/\lambda =
\qty{212}{ka}$. État stationnaire~: $\rho\varepsilon/\Lambda = P_0/N - \lambda =
\qty{4.46e-6}{an^{-1}}$, $\varepsilon = \qty{2.64e-4}{cm.an^{-1}} =
\qty{2.6}{m.Ma^{-1}}$. Pour tout $\varepsilon$ inférieur à cette valeur, $N =
P_0(1 - e^{-\mu t})/\mu$ fixe un $t$, qui va de \qty{212}{ka} à
l'infini~: une mesure a deux inconnues. Un second nucléide de demi-vie
différente, comme l'aluminium-26 dans le même quartz, ou la concentration à
la profondeur dans une carotte, les sépare.
\end{solution}

\begin{solution}{pb:b3:advanced-geochronology:1}
\textbf{1.} $t_{206}$ / $t_{207}$, en \unit{\Ma}~: Z1 301.9 / 302.2, Z2 284.1 /
286.5, Z3 265.7 / 268.5, Z4 246.1 / 250.1, Z5 225.1 / 229.4. Une erreur de
\qty{0.2}{\percent} sur $y$ vaut environ \qty{0.6}{\Ma}~: Z1 est concordant
à l'erreur près~; les autres sont discordants, de plus en plus de Z2 à Z5.

\textbf{2.} Les zircons ont perdu une fraction croissante de leur plomb lors
d'un événement ultérieur~: ils se situent de plus en plus bas sur la corde. La corde se place sous la
\omterm{def:b3:advanced-geochronology:concordia}{concordia} concave, où $x$ est plus grand pour un $y$ donné~: $t_{207} > t_{206}$.

\textbf{3.} Z1~: $X = 1/0.04794 = 20.859$, $Y = 0.3467/(137.8\times0.04794) =
0.05248$. Z5~: $X = 28.145$, $Y = 0.05178$.

\textbf{4.} Poids $1.000$, $1.235$, $1.563$, $1.563$, $2.041$ (en
$10^8$)~; $\bar x = 0.29402$, $\bar y = 0.040892$, $S_{xx} = \num{7.845e5}$.

\textbf{5.} $b = 0.13260$, $a = \bar y - b\bar x = 0.001905$.

\textbf{6.} À \qty{310}{\Ma}, $\mathrm dy/\mathrm dx = 0.1575\,
e^{(\lambda_{238} - \lambda_{235})t} = 0.1218$. La \omterm{def:b3:advanced-geochronology:concordia}{discordia} n'est que
\qty{9}{\percent} plus pentue que la \omterm{def:b3:advanced-geochronology:concordia}{concordia}~: elle la rencontre sous un petit angle,
et l'\omterm{def:b3:advanced-geochronology:concordia}{intercept supérieur} sera imprécis.

\textbf{7.} Newton à partir de \qty{310}{\Ma}~: $t = \qty{311.2}{\Ma}$.

\textbf{8.} À partir de \qty{100}{\Ma}~: $t = \qty{99.7}{\Ma}$.

\textbf{9.} $(y_i - a - bx_i)/\sigma_i = 0.64$, $-0.91$, $0.35$, $-0.24$,
$0.17$~; $\mathrm{MSWD} = 1.435/3 = 0.48$, inférieur à la limite $2.63$ pour $f = 3$~:
les zircons s'alignent sur une seule droite.

\textbf{10.} $s_b = 1/\sqrt{\num{7.845e5}} = 0.00113$.

\textbf{11.} Avec $b - 2s_b = 0.13034$, l'\omterm{def:b3:advanced-geochronology:concordia}{intercept supérieur} vaut
\qty{322.9}{\Ma}~; avec $b + 2s_b = 0.13485$, \qty{302.6}{\Ma}. L'âge est
$311 \pm \qty{10}{\Ma}$ à deux écarts-types.

\textbf{12.} $f = (x(311.2) - x_i)/(x(311.2) - x(99.7))$~: Z1 \qty{4.7}{\percent},
Z5 \qty{41}{\percent}.

\textbf{13.} $J = 0.015754/1.3129 = 0.0119999$, $0.01200$.

\textbf{14.} Étape 1~: \qty{150.0}{\Ma}~; étape 4~: \qty{296.6}{\Ma}~; étape 10~:
\qty{295.0}{\Ma}.

\textbf{15.} Étapes 4 à 10~: sept étapes consécutives, \qty{82}{\percent} de
l'argon-39, sans tendance. Les étapes 1 à 3 montent en escalier.

\textbf{16.} Erreurs égales~: la moyenne vaut $2072.9/7 = \qty{296.1}{\Ma}$, $\sigma =
1.1/\sqrt7 = \qty{0.42}{\Ma}$~; la somme des écarts au carré vaut $3.83$,
$\mathrm{MSWD} = 3.83/(1.21\times6) = 0.53$, inférieur à $2.15$.

\textbf{17.} $\bar R = \sum f_iR_i = 14.270$, $t = \qty{285.2}{\Ma}$~: l'argon
perdu par les bordures, libéré dans les premières étapes, abaisse la moyenne, comme
il abaisserait un âge potassium--argon.

\textbf{18.} $(700 - 375)/(311.2 - 296.1) = \qty{21.5}{\celsius.Ma^{-1}}$
(avec une large fourchette, puisque l'âge du zircon est connu à $\pm\qty{10}{\Ma}$).

\textbf{19.} Vers \qty{100}{\Ma}, le granite était enfoui sous sa couverture
mésozoïque, et un réchauffement ou une circulation de fluides a fait perdre du plomb aux zircons.
La biotite a perdu de l'argon par ses bordures, ce qu'enregistre l'escalier~; sa
plus basse étape, \qty{150}{\Ma}, est un mélange d'argon ancien et remis à zéro. Le plateau
a résisté~: la biotite n'a pas été assez chauffée pour perdre tout son argon.

\textbf{20.} La pente des moindres carrés de la hauteur en fonction de l'âge est
\qty{20.1}{m.Ma^{-1}}.

\textbf{21.} $(65 - 10)/30 = \qty{1.83}{km}$~; depuis \qty{43.1}{\Ma}, une moyenne de
$1830/43.1 = \qty{42}{m.Ma^{-1}}$.

\textbf{22.} $\rho\varepsilon/\Lambda = P_0/N - \lambda = \num{3.415e-5} -
\num{0.050e-5} = \qty{3.365e-5}{an^{-1}}$; $\varepsilon = \num{3.365e-5}\times
160/2.7 = \qty{1.99e-3}{cm.an^{-1}} = \qty{19.9}{m.Ma^{-1}}$.

\textbf{23.} $P_0\Lambda/(\rho N) = \qty{20.2}{m.Ma^{-1}}$, \qty{1.5}{\percent}
de plus~: la désintégration ne joue presque pas à ce rythme.

\textbf{24.} $\Lambda/(\rho\varepsilon) = 160/(2.7\times\num{1.99e-3}) =
\qty{30}{ka}$. Le granite s'est soulevé à environ \qty{20}{m.Ma^{-1}} entre 52 et
\qty{43}{\Ma}, à \qty{42}{m.Ma^{-1}} en moyenne depuis, et s'érode à
\qty{20}{m.Ma^{-1}} aujourd'hui~: la mise à nu a été plus rapide à un moment
du Cénozoïque.

\textbf{25.} Le granite de Varden a cristallisé
il y a \textbf{\qty[text-series-to-math=true, separate-uncertainty=true]{311 \pm 10}{\Ma}} (\omterm{def:b3:advanced-geochronology:concordia}{intercept supérieur},
à deux écarts-types).
\end{solution}
